\chapter{Shortest Paths}

\tikzset{sv/.style={circle, draw, thick, minimum size=7mm, inner sep=0pt, font=\small},
  sw/.style={font=\scriptsize, fill=white, text=black, inner sep=1pt},
  sa/.style={thick, -{Stealth[length=2.2mm]}}}

Let $G=(V,E)$ be a directed graph with edge weights $w$. The weight of a path is the sum of its edge weights, and
$\delta(u,v)$ is the minimum weight of a $u$-$v$ path ($\infty$ if none; $-\infty$ if a negative-weight cycle is reachable on the way).

\begin{center}
\begin{tabular}{llll}
\toprule
Problem & Weights & Algorithm & Time \\
\midrule
Single source & unweighted & BFS & $\bigO(n+m)$ \\
Single source & DAG, any weights & topological order & $\bigO(n+m)$ \\
Single source & $w\ge0$ & Dijkstra & $\bigO((n+m)\log n)$ \\
Single source & any (detects neg.\ cycles) & Bellman--Ford & $\bigO(nm)$ \\
All pairs & any, no neg.\ cycles & Floyd--Warshall & $\Theta(n^3)$ \\
\bottomrule
\end{tabular}
\end{center}

\section{Common structure}
\begin{lemma}[Optimal substructure]
Any subpath of a shortest path is a shortest path.
\end{lemma}
\begin{proof}
If a subpath $x\leadsto y$ could be replaced by a lighter one, the whole path would get lighter.
\end{proof}
If there is no negative cycle, shortest paths can be taken to be simple (at most $n-1$ edges): removing a cycle of nonnegative weight does not increase weight.

All single-source algorithms maintain estimates $d[v]\ge\delta(s,v)$ and improve them by \textbf{relaxation}:
\begin{lstlisting}[style=pseudo]
Relax(u, v):
    if d[u] + w(u, v) < d[v] then
        d[v] = d[u] + w(u, v); parent[v] = u
\end{lstlisting}
Initialize $d[s]=0$, $d[v]=\infty$ otherwise. Key facts: $d[v]$ never drops below $\delta(s,v)$ (it is always the weight of some path), and
(\emph{path-relaxation property}) if $s=v_0,v_1,\dots,v_k$ is a shortest path and its edges are relaxed in this order (possibly interleaved with other relaxations),
then afterwards $d[v_k]=\delta(s,v_k)$.

\section{Shortest paths in a DAG}
\begin{lstlisting}[style=pseudo]
DAGShortestPaths(G, w, s):
    topologically sort G
    initialize d
    for each u in topological order do
        for each edge (u, v) do Relax(u, v)
\end{lstlisting}
Every path visits vertices in topological order, so its edges are relaxed in path order; by path relaxation the result is correct. $\bigO(n+m)$.
Negative weights are fine. Negating weights gives \emph{longest} paths in a DAG (critical path in project scheduling); longest simple path in a general graph is NP-hard.

\section{Dijkstra's algorithm}
For nonnegative weights.
\begin{lstlisting}[style=pseudo]
Dijkstra(G, w, s):
    initialize d; S = {}
    Q = min-priority queue on V keyed by d
    while Q not empty do
        u = ExtractMin(Q); S = S $\cup$ {u}
        for each edge (u, v) do Relax(u, v)     // DecreaseKey if d[v] changes
\end{lstlisting}

\begin{theorem}
If all weights are nonnegative, then when $u$ is extracted, $d[u]=\delta(s,u)$.
\end{theorem}
\begin{proof}
Suppose not, and let $u$ be the first vertex extracted with $d[u]>\delta(s,u)$. Then $u\ne s$ and $u$ is reachable. Take a shortest path $P$ from $s$ to $u$ and let
$(x,y)$ be the first edge of $P$ with $x\in S$, $y\notin S$. When $x$ was extracted, $d[x]=\delta(s,x)$ (choice of $u$) and $(x,y)$ was relaxed, so
$d[y]=\delta(s,y)$ (subpaths of shortest paths are shortest). Since weights are nonnegative, $\delta(s,y)\le\delta(s,u)$. So
\[
d[y]=\delta(s,y)\le\delta(s,u)<d[u],
\]
yet $u$ was extracted before $y$, so $d[u]\le d[y]$. Contradiction.
\end{proof}
Runtime with a binary heap: $n$ extractions and $\le m$ decrease-keys, $\bigO((n+m)\log n)$; with a Fibonacci heap $\bigO(m+n\log n)$; with an array $\bigO(n^2)$.

\begin{figure}[H]
\centering
\begin{tikzpicture}[x=1.6cm, y=1.3cm]
  \node[sv, fill=gray!20] (s) at (0,0.6) {$s$};
  \node[sv] (t) at (1.2,1.5) {$t$}; \node[font=\scriptsize, above=1pt] at (t.north) {8};
  \node[sv] (x) at (2.8,1.5) {$x$}; \node[font=\scriptsize, above=1pt] at (x.north) {9};
  \node[sv] (y) at (1.2,-0.3) {$y$}; \node[font=\scriptsize, below left=0pt] at (y.south west) {5};
  \node[sv] (z) at (2.8,-0.3) {$z$}; \node[font=\scriptsize, below=1pt] at (z.south) {7};
  \node[font=\scriptsize, left=1pt] at (s.west) {0};
  \draw[sa] (s) -- node[sw]{10} (t);
  \draw[sa, very thick, blue!60!black] (s) -- node[sw]{5} (y);
  \draw[sa, very thick, blue!60!black] (t) -- node[sw]{1} (x);
  \draw[sa] (t) to[bend right=15] node[sw]{2} (y);
  \draw[sa, very thick, blue!60!black] (y) to[bend right=15] node[sw]{3} (t);
  \draw[sa] (y) -- node[sw, pos=0.3]{9} (x);
  \draw[sa, very thick, blue!60!black] (y) -- node[sw]{2} (z);
  \draw[sa] (x) to[bend left=15] node[sw]{4} (z);
  \draw[sa] (z) to[bend left=15] node[sw]{6} (x);
  \draw[sa] (z) to[out=-150, in=-80, looseness=1.2] node[sw, pos=0.35]{7} (s);
\end{tikzpicture}
\caption{Dijkstra from $s$ extracts $s, y, z, t, x$ in that order, with
final distances shown next to each vertex. The parent pointers (bold) form
a shortest-path tree.}
\end{figure}

\begin{important}[Negative edges break Dijkstra]
Vertices $s,a,b$ with $w(s,a)=2$, $w(s,b)=3$, $w(b,a)=-2$. Dijkstra extracts $a$ with $d[a]=2$ and never updates it, but $\delta(s,a)=1$.

\begin{center}
\begin{tikzpicture}[x=1.5cm, y=1.1cm]
  \node[sv] (s) at (0,0) {$s$}; \node[sv] (a) at (2,0.6) {$a$}; \node[sv] (b) at (2,-0.6) {$b$};
  \draw[sa] (s) -- node[sw]{2} (a);
  \draw[sa] (s) -- node[sw]{3} (b);
  \draw[sa, red!70!black] (b) -- node[sw]{$-2$} (a);
  \node[font=\scriptsize, anchor=west] at (2.3,0.6) {extracted with $d=2$, but $\delta(s,a)=1$};
\end{tikzpicture}
\end{center}
Adding a constant to every edge does \emph{not} fix this: it penalizes paths with more edges.
\end{important}

\section{Bellman--Ford}
\begin{lstlisting}[style=pseudo]
BellmanFord(G, w, s):
    initialize d
    repeat n - 1 times:
        for each edge (u, v) in E do Relax(u, v)
    for each edge (u, v) in E do
        if d[u] + w(u, v) < d[v] then return "negative cycle reachable from s"
    return d
\end{lstlisting}
Runtime $\Theta(nm)$.

\begin{theorem}
If no negative cycle is reachable from $s$, Bellman--Ford returns $d[v]=\delta(s,v)$ for all $v$; otherwise it reports a negative cycle.
\end{theorem}
\begin{proof}
No negative cycle: each reachable $v$ has a shortest simple path $v_0=s,\dots,v_k=v$ with $k\le n-1$. In round $i$, edge $(v_{i-1},v_i)$ is relaxed, so by path
relaxation $d[v]=\delta(s,v)$ after round $k$. Then no edge can be relaxed further ($d[v]\le d[u]+w(u,v)$ is the triangle inequality for $\delta$), so no cycle is reported.

Negative cycle $v_0,\dots,v_k=v_0$ reachable, and suppose no edge is reported: $d[v_i]\le d[v_{i-1}]+w(v_{i-1},v_i)$ for all $i$. Summing around the cycle,
the $d$ terms cancel (all finite since reachable), giving $0\le\sum w(v_{i-1},v_i)<0$. Contradiction.
\end{proof}

\paragraph{DP view.} Let $D_i[v]$ = minimum weight of an $s$-$v$ path with at most $i$ edges. Then
$D_i[v]=\min\bigl(D_{i-1}[v],\ \min_{(u,v)\in E}D_{i-1}[u]+w(u,v)\bigr)$, and $\delta(s,v)=D_{n-1}[v]$. Bellman--Ford is this DP computed in place.
\emph{Application}: currency arbitrage (weights $-\log(\text{rate})$; an arbitrage is a negative cycle).

\section{All-pairs shortest paths: Floyd--Warshall}
Running Bellman--Ford from every source is $\bigO(n^2m)$, which is $\bigO(n^4)$ for dense graphs; Dijkstra from every source is $\bigO(nm\log n)$ but needs $w\ge0$.

Number the vertices $1..n$. Let $D^{(k)}[i][j]$ be the minimum weight of an $i$-$j$ path whose \emph{intermediate} vertices all lie in $\{1,\dots,k\}$.
A shortest such path either avoids vertex $k$, or passes through $k$ exactly once (no negative cycles), splitting into $i\leadsto k$ and $k\leadsto j$ with intermediates in $\{1..k-1\}$:

\begin{center}
\begin{tikzpicture}[x=1.4cm, y=1cm]
  \node[sv] (i) at (0,0) {$i$}; \node[sv] (j) at (4,0) {$j$}; \node[sv, fill=orange!20] (k) at (2,1.4) {$k$};
  \draw[sa, decorate, decoration={snake, amplitude=1.2pt, segment length=8pt, post length=2mm}] (i) -- node[sw, below=2pt]{$D^{(k-1)}[i][j]$} (j);
  \draw[sa, decorate, decoration={snake, amplitude=1.2pt, segment length=8pt, post length=2mm}] (i) -- node[sw, above left]{$D^{(k-1)}[i][k]$} (k);
  \draw[sa, decorate, decoration={snake, amplitude=1.2pt, segment length=8pt, post length=2mm}] (k) -- node[sw, above right]{$D^{(k-1)}[k][j]$} (j);
  \node[font=\scriptsize, gray] at (2,-0.9) {all wiggly paths use only intermediates in $\{1,\dots,k-1\}$};
\end{tikzpicture}
\end{center}
\[
D^{(0)}[i][j]=\begin{cases}0 & i=j\\ w(i,j) & (i,j)\in E\\ \infty & \text{else}\end{cases},\qquad
D^{(k)}[i][j]=\min\bigl(D^{(k-1)}[i][j],\ D^{(k-1)}[i][k]+D^{(k-1)}[k][j]\bigr).
\]
\begin{lstlisting}[style=pseudo]
FloydWarshall(W):
    D = W
    for k = 1 to n do
        for i = 1 to n do
            for j = 1 to n do
                if D[i][k] + D[k][j] < D[i][j] then
                    D[i][j] = D[i][k] + D[k][j]; next[i][j] = next[i][k]
    return D
\end{lstlisting}
$\Theta(n^3)$ time, $\Theta(n^2)$ space (updating in place is safe since row $k$ and column $k$ do not change in round $k$).
A negative cycle exists iff some $D[i][i]<0$ at the end.
The same recurrence with Boolean operations computes the \textbf{transitive closure} (Warshall).

\begin{aside}
\textbf{Johnson's algorithm} handles sparse graphs with negative edges: add a new vertex $q$ with 0-weight edges to all vertices, run Bellman--Ford from $q$ to get
$h(v)=\delta(q,v)$, reweight $w'(u,v)=w(u,v)+h(u)-h(v)\ge0$ (path weights change by $h(\text{start})-h(\text{end})$, so shortest paths are preserved), then run Dijkstra from every vertex.
Total $\bigO(nm\log n)$.
\end{aside}

\section*{Summary}
\begin{itemize}
  \item All algorithms are built on relaxation; correctness via the path-relaxation property or an exchange/contradiction argument.
  \item DAG: relax in topological order, linear time, any weights.
  \item Dijkstra: greedy, needs nonnegative weights, $\bigO(m\log n)$.
  \item Bellman--Ford: $n-1$ rounds of relaxing all edges, $\bigO(nm)$; extra round detects negative cycles.
  \item Floyd--Warshall: DP over allowed intermediate vertices, $\Theta(n^3)$.
\end{itemize}
